% < < < < < < < < < < < < < < < < < < < < < < < < < < < < < < < < < < < < 
% BEGIN: MULTICOL ENV.
% < < < < < < < < < < < < < < < < < < < < < < < < < < < < < < < < < < < < 

% ---------------------- MULTICOL. ENV. PARAMETERS ----------------------
\begin{multicols*}{3}
\setlength{\columnseprule}{1pt}
\setlength{\premulticols}{1pt}
\setlength{\postmulticols}{1pt}
\setlength{\multicolsep}{1pt}
\setlength{\columnsep}{15pt}
% -----------------------------------------------------------------------

% ---------------------------- SET THE TITLE ----------------------------
\begin{center} % Begin document title
\normalsize{\textsc{\textbf{Computer Science Summary \textendash{} ETH Zürich}}}
\rule[0pt]{\columnwidth}{1.5pt}
\end{center}
\vspace{10pt}
% -----------------------------------------------------------------------

% ----------------------- SET A SECTION TITLE ---------------------------
\section{\textsc{Analysis and Programming}}
\rule[0pt]{\columnwidth}{0.5pt}
\vspace{10pt}
% -----------------------------------------------------------------------

\begin{definition}[Limit]
    A limit is the value that a function approaches as the input approaches some value.
    $\lim_{x \to c} f(x) = L$.
\end{definition}

\begin{theorem}[Taylor Expansion]
    $f(x) = \sum_{n=0}^{\infty} \frac{f^{(n)}(a)}{n!} (x-a)^n$.
\end{theorem}

\begin{lemma}[Squeeze Theorem]
    If $g(x) \leq f(x) \leq h(x)$ and $\lim_{x \to c} g(x) = \lim_{x \to c} h(x) = L$, then $\lim_{x \to c} f(x) = L$.
\end{lemma}

\begin{remark}
    Limit laws apply only when the limits exist.
\end{remark}

\begin{important}
    The Big-O notation describes the limiting behavior of a function.
\end{important}

\begin{example}
    $O(n \log n)$ is the complexity of Merge Sort.
\end{example}

\begin{codeexample}[language=C++]
// Merge Sort Example
void mergeSort(int arr[], int l, int r);
\end{codeexample}

\begin{algorithm}[caption=Binary Search]
Target x in sorted arr[]
1. l=0, r=n-1
2. while l <= r:
   m = (l+r)/2
   if arr[m] == x: return m
   if arr[m] < x: l = m+1
   else: r = m-1
\end{algorithm}

% ---------------------- SET A SUBSECTION TITLE -------------------------
\subsection{\textsc{Complexity Table}}
\hdashrule[0pt]{\columnwidth}{0.5pt}{1mm}
\vspace{10pt}

\begin{center}
    \begin{tabular}{l|c}
    \textbf{Algorithm} & \textbf{Complexity} \\
    \hline\hline
    Quick Sort & $O(n^2)$ / $O(n \log n)$ \\
    Binary Search & $O(\log n)$ \\
    DFS / BFS & $O(V + E)$ \\
    \end{tabular}
\end{center}

\end{multicols*}

% < < < < < < < < < < < < < < < < < < < < < < < < < < < < < < < < < < < < 
% END: MULTICOL ENV.
% < < < < < < < < < < < < < < < < < < < < < < < < < < < < < < < < < < < < 
